\subsection{基域的限制}

本号中给定两个完备的非离散赋值交换域 K 和 L，以及把赋值域 K 同构到 L 的一个子域上的同构 \(\sigma\)。若 E 是 L 上的 Banach 空间，则以 \(\sigma_{*}(\mathbf{E})\) 表示通过限制标量得到的 K-向量空间（参见~代数，第~II~章，第 3 版，§ 8）；E 上给定的拓扑与 K-向量空间结构相容，且 \(\sigma_{*}(\mathbf{E})\) 是 K 上的 Banach 空间。

5.14.1. 设 X 是 L 上的解析流形，\(c = (U, \varphi, E)\) 是 X 的一个图。映射 \(\varphi\) 将 U 双射到 \(\sigma_{*}(E)\) 的一个开子集上，三元组 \(c_{\sigma} = (U, \varphi, \sigma_{*}(E))\) 是 X 的一个图。在 X 上存在唯一的 K 上解析流形结构，使得对于 L-解析流形 X 的每个图 c，\(c_{\sigma}\) 都是图。将如此得到的 K 上解析流形记为 \(X_{\sigma}\)，称 \(X_{\sigma}\) 是通过限制标量（借助 \(\sigma\) 从 L 到 K）由 X 得到的。\(X_{\sigma}\) 所对应的拓扑空间与 X 所对应的拓扑空间相同。

5.14.2. 例：

\begin{enumerate}
\def\labelenumi{(\alph{enumi})}
\item
  取 K = R、L = C，σ 为 R 到 C 的典范嵌入。因此，每个复解析流形都典范地带有实解析流形结构；这个实解析结构本身为每个 r 定义了 \(C^{r}\) 类微分结构。
\item
  取 K = C、L = C，σ 为共轭 \(x \mapsto \bar{x}\)。于是给每个复解析流形 X 关联一个复解析流形 \(\bar{X}\)，称为 X 的共轭。若 f 是定义在 X 的一个开子集 U 上的复值函数，则 f 对 \(\bar{X}\) 的结构解析，当且仅当共轭函数 \(\bar{f}\) 对 X 的结构解析。流形 X 和 \(\bar{X}\) 通过限制标量定义同一个实解析流形。
\end{enumerate}

5.14.3. 设 X 是 K 上的解析流形，Y 是 L 上的解析流形。若映射 \(f:X\to Y\) 是从 X 到  的态射，则称其为 \(\sigma\)-解析的。若 \(Y_{\sigma}\)\(K\subset L\)，σ 是 K 到 L 的嵌入，且 X 是 L 上的解析流形，则称从  到 Y 的 \(\sigma\)-解析映射为 K-解析映射。\(\sigma\)\(X_{\sigma}\)

5.14.4. 设 V 是 L 上的 Banach 空间。有 \(V_{\sigma} = \sigma_{*}(V)\)；\(\sigma_{*}(V)\) 在 K 上的典范解析流形结构由 V 在 L 上的典范解析流形结构通过限制标量得到。

5.14.5. 设 X 是 L 上的解析流形，且 \(a \in X\)。令 c 为 X 在 a 处的一个图；则 \(c_{\sigma}\) 是 \(X_{\sigma}\) 在 a 处的一个图，并由此得到同构

\[
\theta_ {c} \colon \mathrm{E} \rightarrow \mathrm{T} _ {a} (\mathrm{X}), \quad \theta_ {c _ {\sigma}} \colon \sigma_ {*} (\mathrm{E}) \rightarrow \mathrm{T} _ {a} (\mathrm{X} _ {\sigma})
\]

从而得到从 \(\theta_{c_{\sigma}} \circ \sigma_{*}(\theta_{c})^{-1}\) 到 \(\sigma_{*}(\mathrm{T}_{a}(\mathrm{X}))\) 的同构 \(\mathrm{T}_{a}(\mathrm{X}_{\sigma})\)；该同构与 \(c\) 的选择无关；借助记号的滥用，它使我们能够把 \(\mathrm{T}_{a}(\mathrm{X}_{\sigma})\) 与 \(\sigma_{*}(\mathrm{T}_{a}(\mathrm{X}))\)，甚至与 \(\mathrm{T}_{a}(\mathrm{X})\) 等同起来。

若 \(f\) 是从 X 到流形 Y 的 L-解析映射，则 \(f\) 在 \(a\) 处的切线性映射（将 \(f\) 视为从 \(\mathrm{X}_{\sigma}\) 到 \(\mathrm{Y}_{\sigma}\) 的态射）等于 \(\sigma_{*}(\mathrm{T}_{a}(f))\)。

5.14.6. 设 X 和 Y 是 L 上的两个解析流形，且 \(f: \mathbf{X}_{\sigma} \to \mathbf{Y}\) 是连续的 \(\sigma\)-解析映射。假设 K 的特征为 0。则 \(f\) 是 L-解析的充分必要条件是：对于每个 \(a \in \mathbf{X}\)，映射 \(\mathrm{T}_{a}(f)\) 是 L-线性的。

当 K = R、L = C 时（第 5.14.2 号的情形 (a)），有更精确的结果：若 X 和 Y 都是有限维的，且 \(f: X \to Y\) 是 \(C^{1}\) 类映射，其在 X 每一点处的切映射都是 C-线性的，则 \(f\) 是复解析的。

5.14.7. 设 X 是复解析流形，g 是从 X 到自身的映射，满足下列条件：

\begin{enumerate}
\def\labelenumi{(\roman{enumi})}
\item
  有 \(g \circ g = \mathrm{Id}_{\mathbf{X}}\)。
\item
  映射 \(g\) 是从解析流形 \(X\) 到共轭流形 \(\overline{X}\) 的同构（5.14.2）。
\end{enumerate}

X 中满足 \(X_{0}\) 的点 x 组成的集合 \(g(x) = x\)，是 X 所对应的实解析流形的闭解析子流形。对于 \(x \in X_{0}\)，有 \(\mathbf{T}_{x}(\mathbf{X}) = \mathbf{T}_{x}(\mathbf{X}_{0}) \oplus i\mathbf{T}_{x}(\mathbf{X}_{0})\)。

设 U 是 X 的一个连通开子集，f 和 g 是从 U 到复分离局部凸空间或复分离解析流形的两个复解析映射。若 f 和 g 在 \(U \cap X_{0}\) 的一个非空子集上重合，且该子集在 \(X_{0}\) 中是开的，则 f = g。

假设 \(X_{0}\) 是旁紧的。若 f 是从 \(X_{0}\) 到分离局部凸空间或分离复解析流形的实解析映射，则存在 X 中 \(X_{0}\) 的一个开邻域 U，以及从 U 到 f 的值域的复解析映射，延拓 f。任意两个这样的延拓在 X 中 \(X_{0}\) 的某个邻域上重合。

假设 X 是有限维的。对于每个点 \(a \in X_{0}\)，存在 a 的某个开邻域 U 中的一组（复）坐标 \(\zeta^{1}, \ldots, \zeta^{n}\)，使得当 \(a\) 时有 \(\zeta^{i} \circ g = \bar{\zeta}^{i}\)；于是 \(1 \leqslant i \leqslant n\)\(\xi^{i}\)\(\zeta^{i}\) 在 \(U \cap X_{0}\) 上的限制  取实值，而 \((\xi^{1}, \ldots, \xi^{n})\) 是实解析流形 \(a\)\(X_{0}\) 在 a 处的一组坐标。

5.14.8. 对于每个旁紧实解析流形 Y，存在一个有序对 \((X, g)\)，其中 X 是复解析流形，g 是从 X 到自身的映射，满足 5.14.7 的条件 (i) 和 (ii)，并且存在从 Y 到 \(X_{0}\) 的同构 f。称赋予 f 和 g 的 X 为 Y 的复化。

